% GENERATED FILE. Edit canonical lessons, then run npm run generate:books.
% canonical-source-hash: fnv1a64:408b61685c92ea5d
% canonical-lessons: FR-C499-jeu-de-societe, FR-C499-de, FR-C499-pion, FR-C499-echiquier, FR-C499-tricher

\chapter{Jeux de société}
\label{ch:fr-jeux-de-societe}
\hlchaptermodality{499}

\begin{chapteropening}
\textbf{By the end of this chapter:} \emph{I can talk about board games.}

Complete the last lesson of chapter 499: I can talk about board games.
\end{chapteropening}

\section[le jeu de société]{le jeu de société — board game}

\label{lesson:FR-C499-jeu-de-societe}



\begin{quote}
Before the new one: say the French for torch, then the French for campfire.
\end{quote}

\subsection*{You'll want to know: le jeu de société}
\textbf{le jeu de société} — \textquotedblleft{}board game\textquotedblright{}.

\begin{quote}
Le soir, nous jouons à un jeu de société.
\end{quote}

\begin{quote}
\emph{In the evening we play a board game.}
\end{quote}

Plural \emph{des jeux de société}.

\begin{cousinweb}
\emph{le jeu de société} is \textquotedblleft{}a game for company\textquotedblright{}: \emph{jeu} is inherited from Latin \emph{iocus} \textquotedblleft{}a joke, play\textquotedblright{}, and \emph{société} from Latin \emph{societās}, built on \emph{socius} \textquotedblleft{}companion\textquotedblright{} (Lewis \& Short).
\end{cousinweb}

\begin{grammarlens}[title={What you've built}]
One more word for telling, explaining and describing at B1.
\end{grammarlens}

\subsection*{Your turn}
\begin{itemize}
\raggedright
  \item \emph{Say it:} \emph{le jeu de société}
  \item \emph{Say it:} \emph{le jeu de société}, once more
  \item \emph{Say it:} a sentence of your own with \emph{le jeu de société}
  \item \emph{Recall:} say the French for torch, then the French for campfire, then say \emph{le jeu de société} again
\end{itemize}

\begin{tcolorbox}[breakable,colback=teal!4,colframe=teal!35!black,title={Before you move on}]
What does le jeu de société mean? (\textquotedblleft{}Board game\textquotedblright{}.) Say it once more.
\end{tcolorbox}
\section[le dé]{le dé — dice, die}

\label{lesson:FR-C499-de}



\begin{quote}
Before the new one: say the French for campfire, then the French for board game.
\end{quote}

\subsection*{You'll want to know: le dé}
\textbf{le dé} — \textquotedblleft{}dice, die\textquotedblright{}.

\begin{quote}
C'est à toi de lancer le dé.
\end{quote}

\begin{quote}
\emph{It is your turn to throw the dice.}
\end{quote}

Plural \emph{les dés}. \emph{Un dé à coudre} — a thimble.

\begin{cousinweb}
\emph{le dé} is usually traced to Latin \emph{datum} \textquotedblleft{}what is given\textquotedblright{}, from \emph{dare} \textquotedblleft{}to give\textquotedblright{} (Lewis \& Short), but the origin is not settled.
\end{cousinweb}

\begin{grammarlens}[title={What you've built}]
One more word for telling, explaining and describing at B1.
\end{grammarlens}

\subsection*{Your turn}
\begin{itemize}
\raggedright
  \item \emph{Say it:} \emph{le dé}
  \item \emph{Say it:} \emph{le dé}, once more
  \item \emph{Say it:} a sentence of your own with \emph{le dé}
  \item \emph{Recall:} say the French for campfire, then the French for board game, then say \emph{le dé} again
\end{itemize}

\begin{tcolorbox}[breakable,colback=teal!4,colframe=teal!35!black,title={Before you move on}]
What does le dé mean? (\textquotedblleft{}Dice, die\textquotedblright{}.) Say it once more.
\end{tcolorbox}
\section[le pion]{le pion — pawn, counter (in a game)}

\label{lesson:FR-C499-pion}



\begin{quote}
Before the new one: say the French for board game, then the French for dice.
\end{quote}

\subsection*{You'll want to know: le pion}
\textbf{le pion} — \textquotedblleft{}pawn, counter (in a game)\textquotedblright{}.

\begin{quote}
Elle avance son pion de trois cases.
\end{quote}

\begin{quote}
\emph{She moves her counter forward three squares.}
\end{quote}

In school slang, \emph{le pion} is also a supervisor.

\begin{cousinweb}
\emph{le pion} is inherited from Late Latin \emph{pedō} \textquotedblleft{}a foot soldier\textquotedblright{} (EtymDB), built on \emph{pēs} \textquotedblleft{}foot\textquotedblright{} (Lewis \& Short).
\end{cousinweb}

\begin{grammarlens}[title={What you've built}]
One more word for telling, explaining and describing at B1.
\end{grammarlens}

\subsection*{Your turn}
\begin{itemize}
\raggedright
  \item \emph{Say it:} \emph{le pion}
  \item \emph{Say it:} \emph{le pion}, once more
  \item \emph{Say it:} a sentence of your own with \emph{le pion}
  \item \emph{Recall:} say the French for board game, then the French for dice, then say \emph{le pion} again
\end{itemize}

\begin{tcolorbox}[breakable,colback=teal!4,colframe=teal!35!black,title={Before you move on}]
What does le pion mean? (\textquotedblleft{}Pawn, counter (in a game)\textquotedblright{}.) Say it once more.
\end{tcolorbox}
\section[l'échiquier]{l'échiquier — chessboard}

\label{lesson:FR-C499-echiquier}



\begin{quote}
Before the new one: say the French for dice, then the French for pawn.
\end{quote}

\subsection*{You'll want to know: l'échiquier}
\textbf{l'échiquier} — \textquotedblleft{}chessboard\textquotedblright{}.

\begin{quote}
Ils ont posé l'échiquier sur une petite table.
\end{quote}

\begin{quote}
\emph{They put the chessboard on a little table.}
\end{quote}

Masculine. Chess is \emph{les échecs}.

\begin{cousinweb}
\emph{l'échiquier} comes from Old French \emph{eschequier}, built on \emph{eschec} \textquotedblleft{}check\textquotedblright{}, which goes back through Arabic to Persian \emph{šāh} \textquotedblleft{}king\textquotedblright{}.
\end{cousinweb}

\begin{grammarlens}[title={What you've built}]
One more word for telling, explaining and describing at B1.
\end{grammarlens}

\subsection*{Your turn}
\begin{itemize}
\raggedright
  \item \emph{Say it:} \emph{l'échiquier}
  \item \emph{Say it:} \emph{l'échiquier}, once more
  \item \emph{Say it:} a sentence of your own with \emph{l'échiquier}
  \item \emph{Recall:} say the French for dice, then the French for pawn, then say \emph{l'échiquier} again
\end{itemize}

\begin{tcolorbox}[breakable,colback=teal!4,colframe=teal!35!black,title={Before you move on}]
What does l'échiquier mean? (\textquotedblleft{}Chessboard\textquotedblright{}.) Say it once more.
\end{tcolorbox}
\section[tricher]{tricher — to cheat}

\label{lesson:FR-C499-tricher}



\begin{quote}
Before the new one: say the French for pawn, then the French for chessboard.
\end{quote}

\subsection*{You'll want to know: tricher}
\textbf{tricher} — \textquotedblleft{}to cheat\textquotedblright{}.

\begin{quote}
Mon petit frère triche toujours aux cartes.
\end{quote}

\begin{quote}
\emph{My little brother always cheats at cards.}
\end{quote}

\emph{Un tricheur} — a cheat.

\begin{cousinweb}
\emph{tricher} is inherited from Old French \emph{trichier} (EtymDB), usually traced to Vulgar Latin \emph{triccāre}, from Latin \emph{trīcārī} \textquotedblleft{}to make difficulties, to trifle\textquotedblright{} (Lewis \& Short).
\end{cousinweb}

\begin{grammarlens}[title={What you've built}]
One more word for telling, explaining and describing at B1.
\end{grammarlens}

\subsection*{Your turn}
\begin{itemize}
\raggedright
  \item \emph{Say it:} \emph{tricher}
  \item \emph{Say it:} \emph{tricher}, once more
  \item \emph{Say it:} a sentence of your own with \emph{tricher}
  \item \emph{Recall:} say the French for pawn, then the French for chessboard, then say \emph{tricher} again
\end{itemize}

\begin{tcolorbox}[breakable,colback=teal!4,colframe=teal!35!black,title={Before you move on}]
What does tricher mean? (\textquotedblleft{}To cheat\textquotedblright{}.) Say it once more.
\end{tcolorbox}
